% =============================================================================
%  preamble.tex — shared style for the CardioScheduler report
%  Target engine : pdfLaTeX (also compiles under XeLaTeX / LuaLaTeX)
%  Typography    : Classical journal-grade Times Roman via newtx.
%                  newtxtext + newtxmath is the same typographic family
%                  used by IEEE Transactions, Elsevier, Springer LNCS,
%                  and ACM legacy templates. It pairs with Helvetica
%                  (subtle sans accents) and Inconsolata (monospace).
% =============================================================================

% ── Page geometry — classical 1-inch journal margins
\usepackage[a4paper, margin=1in, headheight=14pt]{geometry}

% ── Encoding & language
\usepackage[T1]{fontenc}
\usepackage[utf8]{inputenc}
\usepackage[english]{babel}

% ── Mathematics — must load before newtxmath per the newtx docs.
\usepackage{amsmath}

% ── Typography — Times Roman (newtx) with matching math.
%   `varg` selects the variant Greek glyphs traditionally preferred in
%   mathematical typesetting; the result is the standard scholarly
%   look used across top-tier engineering and CS journals.
\usepackage{newtxtext}
\usepackage[varg]{newtxmath}
\usepackage{bm}                          % bold math
\usepackage[scaled=0.92]{helvet}         % occasional sans accents
% inconsolata.sty is not bundled in every TeX Live image; fall back to
% beramono and finally to Computer Modern typewriter.
\IfFileExists{inconsolata.sty}{\usepackage[scaled=0.95]{inconsolata}}{%
  \IfFileExists{beramono.sty}{\usepackage[scaled=0.85]{beramono}}{%
    \renewcommand{\ttdefault}{cmtt}}}
\linespread{1.10}                        % comfortable, journal-typical
\usepackage{microtype}                   % protrusion + expansion polish
% Allow paragraphs with stubborn unbreakable groups (e.g. inline math,
% \texttt) to add a little extra word-spacing instead of overflowing
% into the margin. Standard journal-style guard.
\setlength{\emergencystretch}{2em}

% ── Colours — disciplined, IEEE-like palette
\usepackage[table]{xcolor}
\definecolor{ruleblue}{RGB}{30, 60, 120}
\definecolor{accent}{RGB}{170, 35, 35}
\definecolor{slate}{RGB}{60, 70, 90}
\definecolor{mutedgray}{RGB}{120, 120, 130}
\definecolor{tablerow}{RGB}{245, 246, 250}
\definecolor{boxbg}{RGB}{248, 248, 250}

% ── Unicode box-drawing for CLI-style diagrams (matching the slides)
\usepackage{pmboxdraw}
\usepackage{fancyvrb}
% Wrap pointer triangles in a fixed-width box so they line up to the
% monospace grid that pmboxdraw uses for ─ │ ┌ ┐ ┴ ┬ etc.
\newcommand{\monosym}[1]{\hbox to 0.5em{\hss\ensuremath{\smash{#1}}\hss}}
\DeclareUnicodeCharacter{25B6}{\monosym{\blacktriangleright}}
\DeclareUnicodeCharacter{25BC}{\monosym{\blacktriangledown}}
\DeclareUnicodeCharacter{25BA}{\monosym{\blacktriangleright}}
\DeclareUnicodeCharacter{25C0}{\monosym{\blacktriangleleft}}

% ── Tables
\usepackage{booktabs}
\usepackage{array}
\usepackage{tabularx}
\usepackage{makecell}
\usepackage{multirow}
\usepackage{colortbl}
\renewcommand{\arraystretch}{1.18}

% ── Figures, captions, layout
\usepackage{graphicx}
\usepackage{float}
\usepackage[font={small},labelfont={bf},textfont={it},
            labelsep=period,format=plain,justification=centerlast,
            singlelinecheck=true]{caption}
\usepackage{subcaption}
\usepackage{wrapfig}

% ── Code listings (used sparingly for pseudocode / config fragments)
%   Muted, B&W-friendly palette in keeping with classical journal style:
%   bold for keywords, italic grey for comments, no garish colour.
\usepackage{listings}
\lstset{
  basicstyle=\ttfamily\footnotesize,
  keywordstyle=\bfseries,
  commentstyle=\itshape\color{mutedgray},
  stringstyle=\color{slate},
  showstringspaces=false,
  breaklines=true,
  frame=single,
  rulecolor=\color{mutedgray!40},
  framesep=4pt,
  backgroundcolor=\color{boxbg},
  numbers=left,
  numberstyle=\tiny\color{mutedgray},
  numbersep=8pt,
  language=C++,
  morekeywords={uint32_t,int64_t,size_t,LexKey,Acuity,Role,Task,Worker,
                Patient,DirectedGraph,EngineConfig,GreedyMatcher,
                FairGreedyMatcher,HungarianMatcher,PriorityOrdering,
                FcfsOrdering,ScheduleEntry,CapacityReport,EngineMetrics,
                Objectives,Status}
}

% ── Algorithm pseudocode
% algorithm2e is not bundled in this TeX Live image. We define a light
% pseudo-code environment built on a tcolorbox + verbatim-friendly font.
% The macros \KwIn, \KwOut, etc. are emulated for compatibility with the
% existing pseudocode blocks.
\newenvironment{pseudocode}[1][]{%
  \begin{tcolorbox}[
    enhanced, breakable,
    colback=white, colframe=ruleblue!70,
    boxrule=0.4pt, arc=2pt,
    fontupper=\small,
    left=10pt, right=10pt, top=4pt, bottom=4pt,
    title={\sffamily\bfseries Algorithm: #1},
    coltitle=white, colbacktitle=ruleblue,
    breakable
  ]\setlength{\parskip}{2pt}\setlength{\parindent}{0pt}\ttfamily}{%
  \end{tcolorbox}}
\newcommand{\KwIn}[1]{\textbf{Input:} #1\\}
\newcommand{\KwOut}[1]{\textbf{Output:} #1\\}
\newcommand{\KwRet}[1]{\textbf{return} #1}

% ── Section heading style — classical serif, journal-style spacing.
\usepackage{titlesec}
\titleformat{\section}
  {\normalfont\large\bfseries}
  {\thesection.}{0.6em}{}
% Star form (\section*) — used by \tableofcontents, \thebibliography, etc.
% Without this, titlesec can silently drop the heading text on some
% TeX Live versions, leaving "Contents" / "References" invisible.
\titleformat{name=\section,numberless}
  {\normalfont\large\bfseries}
  {}{0pt}{}
\titlespacing{\section}{0pt}{1.6ex plus 0.6ex minus 0.2ex}{0.9ex}

\titleformat{\subsection}
  {\normalfont\normalsize\bfseries}
  {\thesubsection}{0.55em}{}
\titlespacing{\subsection}{0pt}{1.1ex plus 0.4ex minus 0.2ex}{0.5ex}

\titleformat{\subsubsection}
  {\normalfont\normalsize\itshape}
  {\thesubsubsection}{0.45em}{}
\titlespacing{\subsubsection}{0pt}{0.8ex plus 0.3ex}{0.3ex}

% ── Headers / footers — light, italic, journal-style.
\usepackage{fancyhdr}
\pagestyle{fancy}
\fancyhf{}
\renewcommand{\headrulewidth}{0.4pt}
\renewcommand{\footrulewidth}{0pt}
\fancyhead[L]{\small\itshape CardioScheduler}
\fancyhead[R]{\small\itshape A Cardiac-Ward Scheduling \& Benchmarking System}
\fancyfoot[C]{\small\thepage}

% ── Hyperlinks — restrained colour palette.
\usepackage{hyperref}
\hypersetup{
  colorlinks=true,
  linkcolor=ruleblue,
  citecolor=ruleblue,
  urlcolor=ruleblue,
  linktoc=all,
  pdftitle={CardioScheduler: A Multi-Strategy Scheduling and Benchmarking
            System for Cardiac-Ward Workflows},
  pdfauthor={Subash C},
  pdfsubject={SL\,203 Project Report},
  pdfkeywords={scheduling, multi-objective, cardiac ward,
               Hungarian algorithm, lex-policy, Pareto}
}

% ── Lists — relaxed enough to read like prose, not a slide bullet stack.
\usepackage{enumitem}
\setlist{itemsep=2pt, topsep=4pt, parsep=0pt, leftmargin=1.6em,
         partopsep=0pt}

% ── Custom macros
\newcommand{\code}[1]{\texttt{\small #1}}
\newcommand{\flag}[1]{\textbf{\textcolor{accent}{#1}}}
\newcommand{\TODO}[1]{\textcolor{accent}{\textbf{[TODO: #1]}}}
\newcommand{\unit}[1]{\,\text{\small #1}}

% ── Custom figure-placeholder environment
% Usage:
%   \figplaceholder{label}{caption}{height}{width-fraction}{description}
% The description appears inside the box; the filename is recorded as a
% LaTeX comment block above each invocation.
% NOTE: the {label} argument is also used as the figure filename; it
% should be a simple identifier with hyphens, not underscores, to avoid
% subscript-mode interpretation when typeset by \texttt{}.
\newcommand{\figplaceholder}[5]{%
  \begin{figure}[H]
    \centering
    \fcolorbox{mutedgray!50}{boxbg}{%
      \parbox[c][#3][c]{#4\textwidth}{%
        \centering
        \vspace*{0.2em}\par
        {\sffamily\bfseries\color{slate} Figure placeholder}\par
        \vspace{0.4em}
        {\small\sffamily #5\par}
        \vspace{0.5em}
        {\scriptsize\sffamily\color{mutedgray}
         figures/\detokenize{#1}.pdf\par}
      }%
    }
    \caption{#2}
    \label{fig:#1}
  \end{figure}%
}

% ── Numbered key-value box (used for "engineering decisions")
\usepackage{tcolorbox}
\tcbuselibrary{breakable, skins}
\newtcolorbox{insightbox}[1][]{
  enhanced, breakable,
  colback=boxbg, colframe=mutedgray!60,
  boxrule=0.4pt, arc=1pt,
  fontupper=\small\itshape,
  left=10pt, right=10pt, top=5pt, bottom=5pt,
  #1
}

% ── Bibliography style
\usepackage[numbers, sort&compress, square]{natbib}
\bibliographystyle{ieeetr}

% =============================================================================
%  end of preamble.tex
% =============================================================================
